\notepage{51}{Central potentials and radial separation}
A spherically symmetric Hamiltonian has the form
\[
 H=\frac{\mathbf P^2}{2m}+V(R),\qquad
 \mathbf P^2\longleftrightarrow-\hbar^2\nabla^2.
\]
The kinetic operator separates into radial and angular parts,
\[
 \mathbf P^2=P_r^2+\frac{L^2}{R^2},\qquad
 P_r^2\longleftrightarrow
 -\hbar^2\frac1{r^2}\frac\partial{\partial r}
       r^2\frac\partial{\partial r}.
\]
Here $P_r^2$ denotes the radial kinetic differential expression.
For a position-space wavefunction,
\begin{align*}
 \mel{r,\theta,\phi}{\mathbf P^2}{\psi}
 &=-\hbar^2\frac1{r^2}\frac\partial{\partial r}
     \left(r^2\frac{\partial\psi}{\partial r}\right)\\
 &\quad-\frac{\hbar^2}{r^2}
   \left[\frac1{\sin\theta}\frac\partial{\partial\theta}
     \left(\sin\theta\frac\partial{\partial\theta}\right)
     +\frac1{\sin^2\theta}\frac{\partial^2}{\partial\phi^2}\right]\psi.
\end{align*}
Thus the angular operator is
\[
 L^2\longleftrightarrow-\hbar^2
 \left[\frac1{\sin\theta}\frac\partial{\partial\theta}
     \left(\sin\theta\frac\partial{\partial\theta}\right)
     +\frac1{\sin^2\theta}\frac{\partial^2}{\partial\phi^2}\right].
\]
The three operators $H,L^2,L_z$ commute, so choose states
$\ket{E,\ell,m}$ with
\[
 \mel{r,\theta,\phi}{L^2}{E,\ell,m}
 =\ell(\ell+1)\hbar^2\braket{r,\theta,\phi}{E,\ell,m}.
\]
Writing $\Psi_{E\ell m}(r,\theta,\phi)
=\braket{r,\theta,\phi}{E,\ell,m}$, the Schr\"odinger equation becomes
\[
 \left[-\frac{\hbar^2}{2m r^2}\frac\partial{\partial r}
       \left(r^2\frac\partial{\partial r}\right)
       +\frac{\ell(\ell+1)\hbar^2}{2mr^2}+V(r)\right]
       \Psi_{E\ell m}=E\Psi_{E\ell m}.
\]
The radial derivative acts only on $r$; the angular part is an
$L^2$ eigenfunction. This motivates separating the radial function from
the spherical harmonic.
